\scanpage{11}
On trouve donc
\[
\boxed{L_E\left(\deleted{\unsure{t|t^{-1}}}\frac{1}{t\unsure{p\rho|p^\rho}}\right)
=\deleted{\unsure{(-t)^n|(-1/t)^n}}\det f_E\,(-t)^n L_E(t)}
\]
[$\rho$ étant le poids de $E$, qui s'introduit via
\[
\check E\simeq E(\rho)\ \bigr]
\]

\noindent\rule{\linewidth}{0.35pt}\par
\centerline{\rule{0.58\linewidth}{0.35pt}}

Soit $E$ un $\unsure{\mathbb{Q}|Q}_{\ell}$-faisceau constructible sur $X$ \deleted{\unsure{compactifiable}}\inserted{propre} sur $\mathbb{F}_p$, alors on aura
\[
\deleted{H_{\unsure{c|{}}}^i(X,DE)\unsure{,|{}}\simeq D(H^{-i}(X,E))}
\]
\margin{par le th. de dualité globale}
\deleted{d'où dans}
\[
\deleted{\chi(DE)=\sum(-1)^iH^i(X,D(E))\simeq D(\chi(E))}
\]
\[
\mathbb{R}f_\ast(DE)=D(\mathbb{R}f_\ast(E))
\]
d'où
\begingroup\csname tagsleft@true\endcsname\[
\left\{\begin{aligned}
&\boxed{L_{DE}(t)=(-t)^{-\chi(E)}\delta(E)\,L_E(t^{-1})}\\
&\chi(E)=\operatorname{rang}\mathbb{R} f_\ast E=\sum(-1)^i\operatorname{rg} H^i(\bar X,\bar E)\\
&\delta(E)=\prod_i\bigl(\text{dét}\ f_{H^i(\bar X,\bar F)}\bigr)^{(-1)^i}=
\varepsilon(E)p^{\sum_i(-1)^i\deleted{\unsure{\rho/2}}\frac{i+\rho}{2}b_i}
\end{aligned}\right.\tag*{\raisebox{1.35\baselineskip}{(1)}}
\]\endgroup
\margin{invariant géométrique}
\margin{invariant arithmétique}
\margin{avec les conjectures \unsure{standart|standard} si $X$ est \deleted{\unsure{projectif|projective}}\inserted{lisse} \unsure{purement de dim $n$} \deleted{\unsure{lisse,}} et $F$ $c^{t}$. \unsure{tordu|tordue} purement de poids $\rho$}
% slanted
\margin{$b_i=\operatorname{rang} H^i(\bar X,\bar F)$}
% archive sheet number: 10

